\section{Introduction}

% in this paragraph we should describe general fairness & bias issues and interpretability problems to position the work.

Machine learning algorithms that optimize for performance (e.g., accuracy) often result in unfair outcomes~\cite{10.1145/3457607}. These algorithms capture biases present in the training datasets causing discrimination toward different groups. As machine learning continues to be adopted into fields where discriminatory treatments can lead to legal penalties, fairness and interpretability have become a necessity and a legal incentive in addition to an ethical responsibility~\cite{barocas2016big,hacker2020explainable}. Existing methods for fair machine learning include applying complex transformations to the data so that resulting representations are fair~\cite{gupta2021controllable,NEURIPS2018_415185ea,roy2019mitigating,advforg, pmlr-v89-song19a}, adding  regularizers to incorporate fairness~\cite{pmlr-v54-zafar17a,kamishima2012fairness,mehrabi2020statistical}, or modifying the outcomes of unfair machine learning algorithms to ensure fairness~\cite{NIPS2016_9d268236}, among others. Here we present an alternative approach \footnote{Code can be found at: \url{https://github.com/Ninarehm/Attribution}}, which works by identifying the significance of different features in causing unfairness and reducing their effect on the outcomes using an attention-based mechanism.
%Note: We should also emphasize something about interpretability aspect; interpretable in the sense that we identify the features whereas other approaches simply do transformation.


% talk about NLP literature and attention mechanism and our approach in relation to that.
With the advancement of transformer models and the attention mechanism \cite{NIPS2017_3f5ee243}, recent research in Natural Language Processing (NLP) has tried to analyze the effects and the interpretability of the attention weights on the decision making process~\cite{wiegreffe-pinter-2019-attention,jain-wallace-2019-attention,serrano-smith-2019-attention,hao2021self}. Taking inspiration from these works, we propose to use an attention-based mechanism to study the fairness of a model. The attention mechanism provides an intuitive way to capture the effect of each attribute on the outcomes. Thus, by introducing the attention mechanism, we can analyze the effect of specific input features on the model's fairness. We form visualizations that explain model outcomes and help us decide which attributes contribute to accuracy vs.\ fairness. We also show and confirm the observed effect of indirect discrimination in previous work \cite{zliobaite2015survey,6175897,ijcai2017-549} in which even with the absence of the sensitive attribute, we can still have an unfair model due to the existence of proxy attributes. Furthermore, we show that in certain scenarios those proxy attributes contribute more to the model unfairness than the sensitive attribute itself.
% maybe we should add some other more interesting finding here)}

% talk about pre-processing appraoch
Based on the above observations, we propose a post-processing bias mitigation technique by diminishing the weights of features most responsible for causing unfairness. We perform studies on datasets with different modalities and show the flexibility of our framework on both tabular and large-scale text data, which is an advantage over existing interpretable non-neural and non-attention-based models. Furthermore, our approach provides a competitive and interpretable baseline compared to several recent fair learning techniques.

% summarize the work
To summarize, the contributions of this work are as follows: (1) We propose a new framework for attention-based classification in tabular data, which is interpretable in the sense that it allows to quantify the effect of each attribute on the outcomes; (2) We then use these attributions to study the effect of different input features on the fairness and accuracy of the models; (3) Using this attribution framework, we propose a post-processing bias mitigation technique that can reduce unfairness and provide competitive accuracy vs. fairness trade-offs; (4) Lastly, we show the versatility of our framework by applying it to large-scale non-tabular data such as text.



% With the advancement of transformer models and the attention mechanism \cite{NIPS2017_3f5ee243}, recent research has tried to analyze the effects of the attention weights on the decision making process and the interpretability of the models \cite{wiegreffe-pinter-2019-attention,jain-wallace-2019-attention,serrano-smith-2019-attention,hao2020self}. However, despite the abundance of research in this area, less attention has been given to the effect of the attention weights into the fairness of a model. It is important to know what contributes to unfairness in decision making processes in order to make fair and informed decisions. In light of this, we aim to study the effect of the attention mechanism on fairness of the models and their contribution to fairness.

% In this work, inspired from the Natural Language Processing (NLP) research, we propose an attention-based framework for classification tasks on tabular data. We then analyze the effect of having an attention mechanism on the fairness of the model which helps us in having more interpretable fair models. In this analysis, we first show and confirm the observed effect of indirect discrimination in previous work \cite{zliobaite2015survey,6175897,ijcai2017-549} in which even with the absence of the sensitive attribute, we can still have an unfair model due to the existence of proxy attributes. We even go further and show that in some scenarios, the proxy contributes more to the unfairness of the model than the sensitive attribute itself.


% We then propose a framework to capture the effect of each attribute to the fairness/unfairness of the model which can help the interpretability of the model with regards to both fairness and accuracy. Based on this framework, we form visualizations that can help us interpret the models along with observations that can help us decide which attributes contribute to accuracy vs fairness. We then use these results to propose a post-processing bias mitigation technique. In addition to tabular data, we also perform studies on text-based data and show the flexibility of our framework on both tabular and large scale text data which is an advantage over existing interpratable non-neural and non-attention based models.


% To summarize, the contributions of this work are as follows: 1. We propose a new framework for attention based classification in tabular data which can give us some form of interpretability on neural based models. 2. We then use this advantage to study the effect of different attributes on the fairness of the models in addition to accuracy. This can give us more interpratable models both based on accuracy and fairness. 3. Based on this advantage, we then create a visualization to be able to interpret models better in the two fairness and accuracy dimensions. 4. Using this interpretable framework, we then propose a post-processing bias mitigation technique that can reduce unfairness and control the accuracy vs fairness trade-offs. 5. Lastly, we show the flexibility of our framework on large scale non-tabular data, such as text, which provides our framework an advantage over existing models that can not handle large scale non-tabular datasets.

